\chapter{Produkte und gespaltene Gruppenerweiterungen}

\section{Direkte Produkte}

\begin{definition}[Äußeres direktes Produkt]
\label{b40guide:def:direct-product}
Für eine mengenindizierte Familie von Gruppen $(G_i)_{i\in I}$
bezeichnet $\prod_{i\in I}G_i$ das kartesische Produkt mit
koordinatenweiser Multiplikation
\[
 (g_i)_{i\in I}(h_i)_{i\in I}=(g_ih_i)_{i\in I}.
\]
Es heißt \emph{äußeres direktes Produkt}. Das Produkt der leeren Familie
ist die triviale Gruppe. Die Projektion auf den $i$-ten Faktor heißt $p_i$.
\end{definition}

\begin{theorem}[Gruppenstruktur des direkten Produkts]
\label{b40guide:thm:product-group}
Das direkte Produkt ist eine Gruppe mit neutralem Element $(1_{G_i})_i$
und Inversen $(g_i)_i^{-1}=(g_i^{-1})_i$.
Jede Projektion $p_i$ ist ein surjektiver Gruppenhomomorphismus.
Die Einbettung eines Faktors in die $i$-te Koordinate, mit neutralen
Elementen in den übrigen Koordinaten, ist ein Gruppenhomomorphismus;
ihr Bild ist ein Normalteiler des Produkts.
\cite{milne}
\end{theorem}

\begin{theorem}[Universelle Eigenschaft des Gruppenprodukts]
\label{b40guide:thm:product-universal}
Zu einer Gruppe $K$ und Homomorphismen $f_i:K\to G_i$ gibt es genau
einen Gruppenhomomorphismus $f:K\to\prod_iG_i$ mit $p_i\circ f=f_i$
für alle $i$. Er ist gegeben durch
\[
 f(k)=(f_i(k))_{i\in I},\qquad
 \ker f=\bigcap_{i\in I}\ker f_i.
\]
Für $I=\varnothing$ ist der rechte Durchschnitt $K$.
Diese Eigenschaft bestimmt das Gruppenprodukt samt Projektionen
bis auf einen eindeutigen mit den Projektionen verträglichen
Gruppenisomorphismus.
\cite{milne}
\end{theorem}

\begin{definition}[Inneres direktes Produkt]
\label{b40guide:def:internal-direct}
Für Untergruppen $N,H\leq G$ heißt $G$ ihr \emph{inneres direktes Produkt},
wenn
\[
 N\times H\longrightarrow G,\qquad (n,h)\longmapsto nh
\]
ein Gruppenisomorphismus ist.
\end{definition}

\begin{theorem}[Kriterium für ein inneres direktes Produkt]
\label{b40guide:thm:internal-direct}
Für $N,H\leq G$ sind äquivalent:
\begin{enumerate}
\item $G$ ist das innere direkte Produkt von $N$ und $H$.
\item $G=NH$, $N\cap H=\{1\}$ und $N,H\trianglelefteq G$.
\item $G=NH$, $N\cap H=\{1\}$ und $nh=hn$ für alle $n\in N$, $h\in H$.
\end{enumerate}
Dann besitzt jedes $g\in G$ eine eindeutige Darstellung $g=nh$, und
\[
 G/N\cong H,\qquad G/H\cong N.
\]
\cite{milne}
\end{theorem}

\section{Semidirekte Produkte}

\begin{definition}[Äußeres semidirektes Produkt]
\label{b40guide:def:semidirect}
Seien $N,H$ Gruppen und $\alpha:H\to\operatorname{Aut}(N)$
ein Gruppenhomomorphismus; schreibe $\alpha_h=\alpha(h)$.
Das \emph{semidirekte Produkt} $N\rtimes_\alpha H$ ist die Menge
$N\times H$ mit
\[
 (n,h)(n',h')=\bigl(n\alpha_h(n'),hh'\bigr).
\]
Die Reihenfolge der Faktoren und diese Multiplikationskonvention
gehören zur Notation.
\end{definition}

\begin{theorem}[Gruppenstruktur und kanonische Faktoren]
\label{b40guide:thm:semidirect-structure}
$N\rtimes_\alpha H$ ist eine Gruppe mit neutralem Element $(1_N,1_H)$
und
\[
 (n,h)^{-1}=\bigl(\alpha_{h^{-1}}(n^{-1}),h^{-1}\bigr).
\]
Die Abbildungen $i(n)=(n,1_H)$ und $s(h)=(1_N,h)$ sind
Gruppeneinbettungen. Die Projektion $p(n,h)=h$ ist ein surjektiver
Homomorphismus mit
\[
 \ker p=i(N),\qquad p\circ s=\operatorname{id}_H,
 \qquad s(h)i(n)s(h)^{-1}=i(\alpha_h(n)).
\]
Die Untergruppen $i(N)$ und $s(H)$ schneiden sich trivial, und jedes
Element hat genau die Form $i(n)s(h)$.
Für die triviale Wirkung $\alpha_h=\operatorname{id}_N$ stimmt diese
Konstruktion mit $N\times H$ überein.
\cite{milne}
\end{theorem}

\begin{theorem}[Homomorphismen aus einem semidirekten Produkt]
\label{b40guide:thm:semidirect-universal}
Seien $a:N\to K$ und $b:H\to K$ Gruppenhomomorphismen mit
\[
 b(h)a(n)b(h)^{-1}=a(\alpha_h(n))
 \qquad(n\in N,\ h\in H).
\]
Dann gibt es genau einen Gruppenhomomorphismus
$u:N\rtimes_\alpha H\to K$ mit $u\circ i=a$ und $u\circ s=b$.
Er ist $u(n,h)=a(n)b(h)$.
Umgekehrt erfüllen die Einschränkungen jedes solchen
Gruppenhomomorphismus diese Verträglichkeitsbedingung.
\cite{milne}
\end{theorem}

\begin{theorem}[Kriterium für ein inneres semidirektes Produkt]
\label{b40guide:thm:internal-semidirect}
Seien $N\trianglelefteq G$ und $H\leq G$. Gilt
$G=NH$ und $N\cap H=\{1\}$, so ist
\[
 \alpha_h(n)=hnh^{-1},\qquad
 N\rtimes_\alpha H\longrightarrow G,\quad(n,h)\longmapsto nh
\]
ein Gruppenisomorphismus. In diesem Fall heißt $H$ ein
\emph{Komplement} zu $N$ und $G$ ein \emph{inneres semidirektes Produkt}.
Die natürliche Abbildung $H\to G/N$ ist ein Gruppenisomorphismus.
Umgekehrt erfüllen die kanonischen Faktoren jedes äußeren semidirekten
Produkts diese Bedingungen.
\cite{milne,conradsemi}
\end{theorem}

\section{Spaltungen und Faserprodukte}

\begin{definition}[Kurze exakte Folge und Spaltung]
\label{b40guide:def:split-exact}
Eine Folge von Gruppenhomomorphismen
\[
 1\longrightarrow N\xrightarrow{i}G\xrightarrow{p}H
 \longrightarrow1
\]
heißt \emph{kurz exakt}, wenn $i$ injektiv, $p$ surjektiv und
$i(N)=\ker p$ ist. Eine \emph{Spaltung} ist ein Gruppenhomomorphismus
$s:H\to G$ mit $p\circ s=\operatorname{id}_H$.
Die Folge heißt \emph{gespalten}, wenn eine solche Spaltung existiert.
\end{definition}

\begin{theorem}[Gespaltene Erweiterungen sind semidirekte Produkte]
\label{b40guide:thm:split-equivalence}
Für eine kurze exakte Folge wie oben sind äquivalent:
\begin{enumerate}
\item Die Folge besitzt eine Spaltung $s$.
\item $\ker p$ besitzt in $G$ ein Komplement $L$.
\item Für einen Homomorphismus $\alpha:H\to\operatorname{Aut}(N)$
existiert ein Gruppenisomorphismus
$\Phi:N\rtimes_\alpha H\to G$ mit
$\Phi(n,1_H)=i(n)$ und $p(\Phi(n,h))=h$.
\end{enumerate}
Für eine gewählte Spaltung erhält man $L=s(H)$ und
\[
 \alpha_h(n)=i^{-1}\bigl(s(h)i(n)s(h)^{-1}\bigr),\qquad
 \Phi(n,h)=i(n)s(h).
\]
Hier bezeichnet $i^{-1}$ die Umkehrabbildung von $i:N\to i(N)$.
Die Wirkung $\alpha$ ist dabei an die gewählte Spaltung gebunden.
\cite{milne,conradsplit}
\end{theorem}

\begin{definition}[Faserprodukt von Gruppen]
\label{b40guide:def:pullback}
Für Gruppenhomomorphismen $f:G\to Q$ und $h:H\to Q$ heißt
\[
 G\times_Q H=\{(g,a)\in G\times H:f(g)=h(a)\}
\]
mit der vom direkten Produkt übernommenen Multiplikation
\emph{Faserprodukt}. Seine Projektionen werden mit $\pi_G$ und $\pi_H$
bezeichnet.
\end{definition}

\begin{theorem}[Gruppenstruktur und universelle Eigenschaft des Faserprodukts]
\label{b40guide:thm:pullback-universal}
$G\times_QH$ ist eine Untergruppe von $G\times H$, und
$f\circ\pi_G=h\circ\pi_H$.
Für jede Gruppe $K$ und Homomorphismen $a:K\to G$, $b:K\to H$
mit $f\circ a=h\circ b$ gibt es genau einen Gruppenhomomorphismus
\[
 u:K\longrightarrow G\times_QH,\qquad u(k)=(a(k),b(k)),
\]
mit $\pi_G\circ u=a$ und $\pi_H\circ u=b$.
Ist $f$ surjektiv, so ist
\[
 1\longrightarrow\ker f\xrightarrow{\,n\mapsto(n,1_H)\,}
 G\times_QH\xrightarrow{\pi_H}H\longrightarrow1
\]
eine kurze exakte Folge.
\cite{milne}
\end{theorem}